\documentclass[11pt]{article}

\usepackage[a4paper,margin=1in]{geometry}
\usepackage{lmodern}
\usepackage{microtype}
\usepackage{amssymb}
\usepackage{amsmath}
\usepackage{xcolor}
\usepackage{listings}
\usepackage[most]{tcolorbox}
\usepackage{enumitem}
\usepackage{booktabs}
\usepackage{tabularx}
\usepackage{array}
\usepackage{hyperref}

\setlength{\emergencystretch}{3em}

\hypersetup{
  colorlinks=true,
  linkcolor=black,
  urlcolor=teal,
  pdftitle={The Swarm},
}

% ---- colours ---------------------------------------------------------
\definecolor{codebg}{rgb}{0.97,0.97,0.95}
\definecolor{kw}{rgb}{0.10,0.20,0.65}
\definecolor{str}{rgb}{0.15,0.45,0.20}
\definecolor{com}{rgb}{0.45,0.45,0.45}
\definecolor{rule}{rgb}{0.75,0.75,0.72}

% ---- code listings ---------------------------------------------------
\lstdefinelanguage{TS}{
  morekeywords={async,await,function,const,let,var,return,for,of,new,class,
    interface,type,import,export,from,if,else,try,catch,throw,Promise,void,
    switch,case,default,readonly,break,while,true,false,null},
  sensitive=true,
  morecomment=[l]{//},
  morecomment=[s]{/*}{*/},
  morestring=[b]",
  morestring=[b]',
  morestring=[b]`,
}
\lstdefinestyle{mono}{
  backgroundcolor=\color{codebg},
  basicstyle=\ttfamily\small,
  keywordstyle=\color{kw}\bfseries,
  stringstyle=\color{str},
  commentstyle=\color{com}\itshape,
  showstringspaces=false,
  columns=fullflexible,
  keepspaces=true,
  frame=single,
  rulecolor=\color{rule},
  framesep=6pt,
  xleftmargin=10pt,
  xrightmargin=4pt,
  breaklines=true,
  aboveskip=1em,
  belowskip=1em,
}
\lstset{style=mono}

% ---- callout box -----------------------------------------------------
\newtcolorbox{keybox}[1][]{
  colback=codebg, colframe=rule, boxrule=0.5pt, arc=2pt,
  left=8pt,right=8pt,top=6pt,bottom=6pt,
  fonttitle=\bfseries, title={#1}
}

\newcommand{\code}[1]{\texttt{#1}}

\title{\textbf{The Swarm}\\[4pt]
  \large An LLM captain commanding many agents under fog of war}
\author{}
\date{\today}

\begin{document}
\maketitle
\thispagestyle{empty}

\begin{abstract}
\noindent The scout/fighter fleet proved delegation and amalgamation under partial
observability, but its captain's choices were \emph{forced}: one scout to the nearest
frontier, one fighter to the one blocking monster, march when the exit is reachable.
Nothing was decided, so making that captain a model would have been theatre. The
\emph{swarm} removes the forcing. Five scouts, two fighters, and a captain explore a
large branching maze in which the captain faces genuine \emph{allocation under partial
information}: which scout to send where, which fighter to which newly discovered
monster, and when to stop exploring and commit. The captain is an LLM. It sits at the
\emph{orchestration} layer, not the leaf, commanding a real delegation hierarchy in
which it \emph{cannot} cheat, because every order it gives is decoded against the fog
it has been told and bounded by the same reason-gate as a leaf. The construction
sharpens three claims. The amalgamation becomes a true \emph{N-ary} join, and the
captain's map is well-defined precisely because that join is commutative, associative
and idempotent. The scheduler is a tick loop whose determinism rests on a fixed agent
order and serial execution. And the model-as-oracle-behind-a-decoder argument, which
previously guarded a single leaf, now guards the commander: a hallucinated order is a
refused position, never an escape hatch. Everything scripted is deterministic ---
\code{verify.sh} replays it to the same bytes, website included --- and the whole
thing has been run live against a free model.
\end{abstract}

\section{The one idea}

A container is a pair $(P, R)$: a set of prompts $P$ and, for each $p$, a type $R\,p$
of permitted responses. A program is a dependent function $(p : P) \to R\,p$. An
\emph{open agent} is a morphism of containers: a map $u : P \to Q$ sending a prompt
down to a sub-prompt (control flows \emph{down}) and a map $f : (p:P) \to T(u\,p) \to
R\,p$ sending the sub-answer back up (data flows \emph{up}, contravariantly). A
captain delegating to officers, an interpreter calling a device driver, an LLM calling
a tool: one construction.

The fleet put a \emph{scout} at one leaf of this hierarchy and asked whether a model
could make the scout's junction choices. It could. The swarm asks the sharper
question: can a model sit at the \emph{root} --- the commander that delegates to the
whole hierarchy --- and still be bound by the type discipline rather than escaping it?
The answer is the subject of this note. The load-bearing fact is that a model is an
untyped $\mathrm{String} \to \mathrm{String}$; presenting it as a typed planner
$\mathrm{plan} : S \to P\,s$ is a \emph{promise}, and the promise is kept by a
validator. A reply that does not decode throws; it never defaults into a
checked-looking value.

\section{What the swarm adds over the fleet}

The swarm reuses the fleet's world, fog, pathing, two-layer log and combinators
wholesale. Four things are new.

\begin{center}
\begin{tabularx}{\linewidth}{@{}>{\bfseries}l X@{}}
\toprule
Contention & Five scouts over one fog: the captain must fan them to \emph{different}
regions. Two fighters over several monsters: it must allocate, and choose which
monsters are worth a fighter at all.\\
\midrule
N-ary join & The captain folds \emph{all five} scout reports into its fog each step.
The fleet merged one at a time and never exercised associativity; the swarm folds over
$N$, and the fold is well-defined only because the join's laws hold.\\
\midrule
Tick scheduler & The fleet's two-phase round (one scout to completion, then one
fighter, then march) does not survive contention. A pool of eight agents is advanced
one tick at a time, and the captain re-plans around whoever has gone idle.\\
\midrule
Model at the root & The captain's decision is a seam. The default is a deterministic
greedy planner; \code{-{}-live} swaps in a model. The decoder and reason-gate that
guarded the scout now guard the commander.\\
\bottomrule
\end{tabularx}
\end{center}

\section{The N-ary amalgamation}

The captain's entire state is a \code{FogMap}: a flat semilattice of terrain facts
($\mathit{unknown} \sqsubseteq \{\mathit{floor},\mathit{wall},\mathit{exit}\}$, with a
top $\mathit{conflict}$ so the join is total) lifted cellwise, plus a revisable
dictionary of monster positions. The join \code{mergeFog} is idempotent, commutative,
associative and monotone --- law-tested in \code{fog\_test.ts}. This is the
treasury/org-chart ``real algebra'' pole, not a tuple gather.

The fleet merged one scout report at a time, so it only ever used the join as a binary
operation and never relied on associativity. The swarm gathers five reports in a single
planning step and folds them all in:
\[
  \mathrm{fog}' \;=\; \mathrm{fog} \;\sqcup\; \bigsqcup_{i=1}^{5} \mathrm{fogOf}(r_i),
\]
implemented as \code{reports.reduce((f, r) => mergeFog(f, fogOf(r)), base)}.
Because $\sqcup$ is commutative and associative, the result does not depend on the
order the reports arrive, and because it is idempotent, re-merging a report is
harmless. This is not a nicety; it is what \emph{licenses} the scheduler to process the
scouts in a fixed serial order (next section) without the captain's map depending on
that order. The order-independence is proven end-to-end on sightings taken off the
real maze in \code{amalgamate\_test.ts}: folding six reports in natural, reversed and
shuffled orders yields the identical fog.

A subtlety worth stating: within one tick each scout navigates on the fog
\emph{snapshotted at the start of the tick} plus its own new sightings --- it does
\emph{not} see the other scouts' sightings from the same tick. So the five reports are
genuinely independent (a tensor, $\otimes$, not a product), and the join combines them
afterward. Kills are applied outside the fold (a monster the fighter clears is
retracted from the dictionary), keeping the fold itself a pure monotone climb.

\section{The tick scheduler and the determinism contract}

Each agent carries a runtime state --- \code{idle}, \code{exploring},
\code{travelling}, \code{fighting}, \code{done} --- and, for scouts and fighters, a
standing order. One tick advances every active agent by its speed, folds the scouts'
reports into the fog, applies any kills, and --- when there is something to decide ---
asks the captain to re-plan.

\begin{keybox}[The rule the whole thing rests on]
Agents are processed in a \textbf{fixed id order} within a tick, and their gated CLI
calls run \textbf{serially}. \code{tensorC} means ``all five scout reports are owed,''
not ``executed concurrently.'' A \code{Promise.all} over agents that mutate the shared
world file would make the move order --- and therefore the audit --- non-reproducible.
The join's order-independence says the \emph{fog} would be the same; the \emph{log}
would not. So we serialise.
\end{keybox}

One planning step reads as $\mathrm{seqC}(\mathrm{tensorC}(\text{scout}_1 \ldots
\text{scout}_5), \text{dispatch})$: gather all five reports (the serial tensor), fold
them into the fog (the join), then let the captain choose the fighter dispatches and
the march (the dependent \code{next}). The captain commands \emph{only idle agents} and
never reassigns one mid-order, so re-planning can happen as often as we like without
thrashing a standing order.

\paragraph{Event-gated planning.} Re-planning every tick is cheap for the scripted
captain but would mean a model call per tick on \code{-{}-live}. The scheduler therefore
consults the captain only when there is genuinely something to decide. The gate oracle
is the \emph{greedy} plan itself: if the keyless, deterministic scripted captain finds
nothing to assign --- no idle scout with an open frontier, no idle fighter with a
reachable monster, no safe march --- then an idle fighter held in reserve is not an
event, and the captain is not consulted. This bounds the LLM to a couple of dozen calls
over a solve rather than a couple of hundred, and it is \emph{byte-identical} for the
scripted captain: on a skipped tick its plan was empty and changed no move, no log and
no frame.

\section{The scripted captain}

The default \code{CaptainBrain} is a greedy planner, deterministic and keyless, so
\code{verify.sh} stays byte-identical. It operates in two modes.

\textbf{Explore.} Fan idle scouts across \emph{open} frontiers --- those reachable
without a fight, assigned nearest-first so the scouts spread. While open frontier
remains, spend no fighter: exploring is free. Only when exploration \emph{stalls}
(every reachable frontier sits behind a monster) dispatch the nearest idle fighter to
the nearest reachable monster, to reopen the map.

\textbf{Commit.} Once a known-floor corridor to the exit exists, stop exploring and
clear only the monsters \emph{on that corridor}; ignore every other monster. March the
instant a monster-free route is known.

The honest limitation is worth stating plainly. From the fog alone, a monster that
gates the half holding the exit and a monster that gates the decoy half are
\emph{indistinguishable}: both have unknown cells behind them, because the monster is
exactly \emph{why} that ground is dark. So the greedy captain cannot pre-judge which
gate to skip before clearing one --- it must open both, which is precisely what puts
\emph{both fighters} to work at once. What it \emph{can} do honestly is commit-mode
prioritisation: once a monster-free route to the exit is known, it leaves every
off-route monster alone --- the decoy half and the monsters deep in it. Sharper
judgement --- inferring which gate is the decoy before clearing it, and holding a
fighter back --- is where a model has room to do better, and room to do worse.

\medskip\noindent
This is why the shipped maze forks at a central start into two large halves rather
than running a single corridor. An earlier draft placed the monsters on one linear
path to the exit; in a perfect maze the route to the second passes through the first,
so the monsters surfaced strictly one at a time, a single fighter cleared them all, and
the second fighter never moved --- the allocation the swarm is built to show never
arose. The fork forces it: two gates at equal shallow depth, discovered together,
neither distinguishable from fog, so the captain allocates both fighters in one step.

\section{The LLM captain}

\code{modelCaptainBrain} is a second \code{CaptainBrain} behind the same seam; the
scheduler cannot tell it from the greedy default. It reuses the fleet's model stack ---
the neutral \code{ModelClient} boundary and the free, OpenAI-compatible Groq backend
with its 429 backoff --- verbatim.

Each step the captain's fog and roster are serialised into a prompt (an ASCII map with
the agents, monsters and frontier marked, plus explicit lists of idle agents, known
monsters and some frontier coordinates), and the model is offered four tools:
\code{assign\_scout}, \code{dispatch\_fighter}, \code{march}, \code{hold}. The set is
the permission grant; there is no ``do anything.'' Every item of the reply is then
\emph{decoded against the real state}:

\begin{itemize}[leftmargin=1.4em,itemsep=2pt]
  \item a scout id that is not an idle scout --- dropped;
  \item a target cell that is not reachable known floor --- dropped;
  \item a fighter id that is not an idle fighter, or a monster id not known-alive ---
    dropped;
  \item a \code{march} when no monster-free route is known --- dropped;
  \item any order with an empty reason --- dropped (the reason-gate, on orders).
\end{itemize}

An order that does not decode is a refused position, never an escape hatch: the model
can mis-prioritise, but it cannot name a scout that does not exist, send a unit into the
dark, or march through a wall.

\paragraph{Anti-stall, not anti-idle.} The fallback is the delicate part. If the greedy
plan always filled every slot the model left empty, the model could never \emph{choose}
to leave a fighter in reserve --- its prioritisation would be silently overridden and it
would behave exactly like greedy. So the rule is: honour the model's valid orders
as-is, including the choice to leave agents idle; fall back to greedy \emph{only when
the model's whole plan is unproductive while real work remains}. A mis-prioritising
model is slower but honoured; a lazy or stuck model can never deadlock the swarm. This
is the orchestration-layer analogue of the scout's ``off-menu reply $\to$ BFS
fallback,'' raised one level and with the stakes raised to match.

\section{The two-layer log}

Nothing about the audit changes, which is the point. Every move and attack by all five
scouts, both fighters and the captain passes the registry allowlist and the reason-gate,
writes a committed effect to Layer 1 (the ship's log, \emph{what}) and a write-ahead
intent record to Layer 2 (the intent-log, \emph{why}), joined by a correlation id.
Reconciliation cross-checks the two: a committed mutation with no logged reason is an
alarm that must never occur, and a refused mutation is visible only in Layer 2 ---
information the committed log alone could never show. On \code{-{}-live} the id and clock
seams are the \emph{real} ones, never the fixed stand-ins: a live audit must carry true
timestamps, or it is a false record of the very thing reconciliation exists to check.

\section{Live validation}

Run against Groq's \code{openai/gpt-oss-120b}, the LLM captain fanned the five scouts to
different frontiers --- its own reasons, e.g.\ \emph{``monster blocks potential route to
exit''} --- and then, in a single planning step, dispatched \code{fighter-1} to one gate
and \code{fighter-2} to the other: \textbf{both fighters allocated at once}, the decision
the fork exists to force. It cleared the two gates and the one monster on the exit route,
left the decoy monsters alone, and reached the exit in \textbf{136 ticks across 18
planning calls}, with \textbf{493 of 1189 cells never revealed} to it and reconciliation
clean (\textbf{0 alarms}) over a real audit trail. The model commanded as well as the
greedy captain here; the claim the construction actually makes is stronger and humbler at
once: it \emph{could not} command worse than refused.

\section{What is, and is not, decided}

The swarm's thesis is that a language model can run a real delegation hierarchy in which
the type discipline, not trust, is what bounds it --- and that the bound holds whether
the model is strong, weak or adversarial, because a bad order is a refused position. That
is demonstrated. What is \emph{not} claimed is that the model orchestrates \emph{better}
than the greedy planner; on this maze the two land in the same place, and ``the model
decides better'' is a separate, harder, model-dependent question the guardrails
deliberately do not need answered. The honest shape of the result is the same as the
scout's: a free model can slow a solve or mis-prioritise within it, but it never breaks
the solve and never escapes the audit.

\end{document}
